\documentclass{article}

% packages
\usepackage{amsmath, amsthm, thmtools, amsfonts, amssymb, luacode, catchfile, hyperref, ifthen, graphicx, bbm}
\ifcsname c@kobocompile\endcsname
	\usepackage[a5paper, total={1072pt, 1448pt}, margin=10pt, includeheadfoot]{geometry} % set page margins
\else
	\usepackage[a4paper, margin=50pt, includeheadfoot]{geometry}
\fi
\usepackage[shortlabels]{enumitem}
\usepackage[skip=3pt, indent=0pt]{parskip}

% language
\usepackage[bidi=basic, layout=tabular, provide=*]{babel}
\ifcsname c@english\endcsname
	\babelprovide[main, import]{english}
\else
	\babelprovide[main, import, alph=letters.gershayim]{hebrew}
	\setlist[enumerate,2]{label={\alph*.}}
	\babelprovide{rl}
\fi
\babelfont{rm}[Renderer=Harfbuzz]{Libertinus Serif}
\babelfont{sf}{Libertinus Sans}
\babelfont{tt}{Libertinus Mono}

% style
\ifcsname c@stoppagebreaks\endcsname
\else
	\AddToHook{cmd/section/before}{\clearpage}	% Add line break before section
\fi
\linespread{1.3}
\setcounter{secnumdepth}{0}		% Remove default number tags from sections, this won't do well with theorems
\AtBeginDocument{\setlength{\belowdisplayskip}{3pt}}
\AtBeginDocument{\setlength{\abovedisplayskip}{3pt}}
\graphicspath{ {../images/} }

% operators
\DeclareMathOperator\cis{cis}
\DeclareMathOperator\Sp{Span}
\DeclareMathOperator\tr{tr}
\DeclareMathOperator\im{Im}
\DeclareMathOperator\re{Re}
\DeclareMathOperator\diag{diag}
\DeclareMathOperator*\lowlim{\underline{lim}}
\DeclareMathOperator*\uplim{\overline{lim}}
\DeclareMathOperator\rng{rng}
\DeclareMathOperator\Sym{Sym}
\DeclareMathOperator\Arg{Arg}
\DeclareMathOperator\Log{Log}
\DeclareMathOperator\dom{dom}
\DeclareMathOperator\supp{Supp}
\DeclareMathOperator\var{Var}
\DeclareMathOperator\cov{Cov}
\DeclareMathOperator\End{End}
\DeclareMathOperator\id{id}
\DeclareMathOperator\cl{cl}
\DeclareMathOperator\gal{Gal}
\DeclareMathOperator\dist{dist}
\DeclareMathOperator\form{form}
\DeclareMathOperator\term{term}
\DeclareMathOperator\sent{sent}
\DeclareMathOperator\Th{Th}
\DeclareMathOperator\Hom{Hom}
\DeclareMathOperator\Aut{Aut}
\DeclareMathOperator\Char{char}
\DeclareMathOperator\vol{vol}
\DeclareMathOperator\GL{GL}
\DeclareMathOperator\SL{SL}
\DeclareMathOperator\lcm{lcm}

% commands
\newcommand{\cupdot}{\mathbin{\mathaccent\cdot\cup}}
\newcommand{\indicator}{\mathbbm{1}}
\newcommand{\AAA}[0]{\mathbb{A}}
\newcommand{\Aa}[0]{\mathcal{A}}
\newcommand{\BB}[0]{\mathbb{B}}
\newcommand{\Bb}[0]{\mathcal{B}}
\newcommand{\CC}[0]{\mathbb{C}}
\newcommand{\Cc}[0]{\mathcal{C}}
\newcommand{\DD}[0]{\mathbb{D}}
\newcommand{\Dd}[0]{\mathcal{D}}
\newcommand{\EE}[0]{\mathbb{E}}
\newcommand{\FF}[0]{\mathbb{F}}
\newcommand{\Ff}[0]{\mathcal{F}}
\newcommand{\Ii}[0]{\mathcal{I}}
\newcommand{\Gg}[0]{\mathcal{G}}
\newcommand{\HH}[0]{\mathbb{H}}
\newcommand{\Hh}[0]{\mathcal{H}}
\newcommand{\Kk}[0]{\mathcal{K}}
\newcommand{\Ll}[0]{\mathcal{L}}
\newcommand{\Mm}[0]{\mathcal{M}}
\newcommand{\NN}[0]{\mathbb{N}}
\newcommand{\Nn}[0]{\mathcal{N}}
\newcommand{\PP}[0]{\mathbb{P}}
\newcommand{\Pp}[0]{\mathcal{P}}
\newcommand{\QQ}[0]{\mathbb{Q}}
\newcommand{\RR}[0]{\mathbb{R}}
\newcommand{\Rr}[0]{\mathcal{R}}
\newcommand{\Ss}[0]{\mathcal{S}}
\newcommand{\TT}[0]{\mathbb{T}}
\newcommand{\Tt}[0]{\mathcal{T}}
\newcommand{\UU}[0]{\mathbb{U}}
\newcommand{\Uu}[0]{\mathcal{U}}
\newcommand{\Vv}[0]{\mathcal{V}}
\newcommand{\Ww}[0]{\mathcal{W}}
\newcommand{\Xx}[0]{\mathcal{X}}
\newcommand{\Yy}[0]{\mathcal{Y}}
\newcommand{\ZZ}[0]{\mathbb{Z}}
\newcommand{\Zz}[0]{\mathcal{Z}}
\newcommand{\acts}[0]{\circlearrowright}
\newcommand{\explain}[2] {
	\begin{flalign*}
		 && \text{#2} && \text{#1}
	\end{flalign*}
}
\newcommand{\maketitleprint}[1][green]{
	\begin{center}
		\includegraphics[width=6cm]{dragon_#1}
	\end{center}
}

% theorem commands
\newtheoremstyle{c_remark}
	{}	% Space above
	{}	% Space below
	{}% Body font
	{}	% Indent amount
	{\bfseries}	% Theorem head font
	{}	% Punctuation after theorem head
	{.5em}	% Space after theorem head
	{\thmname{#1}\thmnumber{ #2}\thmnote{ \normalfont{\text{(#3)}}}}	% head content
\newtheoremstyle{c_definition}
	{3pt}	% Space above
	{3pt}	% Space below
	{}% Body font
	{}	% Indent amount
	{\bfseries}	% Theorem head font
	{}	% Punctuation after theorem head
	{.5em}	% Space after theorem head
	{\thmname{#1}\thmnumber{ #2}\thmnote{ \normalfont{\text{(#3)}}}}	% head content
\newtheoremstyle{c_plain}
	{3pt}	% Space above
	{3pt}	% Space below
	{\itshape}% Body font
	{}	% Indent amount
	{\bfseries}	% Theorem head font
	{}	% Punctuation after theorem head
	{.5em}	% Space after theorem head
	{\thmname{#1}\thmnumber{ #2}\thmnote{ \text{(#3)}}}	% head content

\ifcsname c@english\endcsname
	\theoremstyle{plain}
	\newtheorem{theorem}{Theorem}[section]
	\newtheorem*{theorem*}{Theorem}
	\newtheorem{lemma}[theorem]{Lemma}
	\newtheorem{proposition}[theorem]{Proposition}
	\newtheorem*{proposition*}{Proposition}

	\theoremstyle{definition}
	\newtheorem{definition}[theorem]{Definition}
	\newtheorem*{definition*}{Definition}
	\newtheorem{example}{Example}[section]
	\newtheorem{exercise}{Exercise}[section]

	\theoremstyle{remark}
	\newtheorem*{remark}{Remark}
	\newtheorem*{solution}{Solution}
	\newtheorem{corollary}[theorem]{Corollary}
	\newtheorem{notation}[theorem]{Notation}
\else
	\theoremstyle{c_plain}
	\newtheorem{theorem}{משפט}[section]
	\newtheorem*{theorem*}{משפט}
	\newtheorem{lemma}[theorem]{למה}
	\newtheorem{proposition}[theorem]{טענה}
	\newtheorem*{proposition*}{טענה}

	\theoremstyle{c_definition}
	\newtheorem{definition}[theorem]{הגדרה}
	\newtheorem*{definition*}{הגדרה}
	\newtheorem{example}{דוגמה}[section]
	\newtheorem{exercise}{תרגיל}[section]

	\theoremstyle{c_remark}
	\newtheorem*{remark}{הערה}
	\newtheorem*{solution}{פתרון}
	\newtheorem{corollary}[theorem]{מסקנה}
	\newtheorem{notation}[theorem]{סימון}
\fi

% Questions related commands
\newcounter{question}
\setcounter{question}{1}
\newcounter{sub_question}
\setcounter{sub_question}{1}

\ifcsname c@english\endcsname
	\newcommand{\question}[1][0]{
		\ifthenelse{#1 = 0}{}{\setcounter{question}{#1}}
		\section{Question \arabic{question}}
		\addtocounter{question}{1}
		\setcounter{sub_question}{1}
	}

	\newcommand{\subquestion}[1][0]{
		\ifthenelse{#1 = 0}{}{\setcounter{sub_question}{#1}}
		\subsection{Part \alph{sub_question}}
		\addtocounter{sub_question}{1}
	}
\else
	\newcommand{\question}[1][0]{
		\ifthenelse{#1 = 0}{}{\setcounter{question}{#1}}
		\section{שאלה \arabic{question}}
		\addtocounter{question}{1}
		\setcounter{sub_question}{1}
	}

	\newcommand{\subquestion}[1][0]{
		\ifthenelse{#1 = 0}{}{\setcounter{sub_question}{#1}}
		\subsection{סעיף \localecounter{letters.gershayim}{sub_question}}
		\addtocounter{sub_question}{1}
	}
\fi

% import lua and start of document
\directlua{common = require ('../common')}

\GetEnv{AUTHOR}

% headers
\author{\AUTHOR}
\date\today

\title{פתרון מטלה 06 --- יסודות באלגברה קומוטטיבית}

\DeclareMathOperator{\loc}{loc}
\DeclareMathOperator{\Frac}{Frac}
\DeclareMathOperator{\Ann}{Ann}

\begin{document}
\maketitle
\maketitleprint[red]

\question{}
יהיו $S, T \subset R$ כפליות, ונסמן $\loc_S$ את העתקת המיקום עבור $S$, וכן $U = \loc_S(T)$.
נראה שמתקיים,
\[
	U^{-1}(S^{-1} R)
	\cong {(S T)}^{-1} R
\]
\begin{proof}
	נגדיר את ההעתקה,
	\[
		\varphi : U^{-1}(S^{-1} R) \to {(S T)}^{-1} R,
		\quad
		\varphi(\frac{r}{s} / \frac{t}{1})
		= \frac{r}{s t}
	\]
	עלינו להראות שההעתקה מוגדרת היטב, וכן שהיא הומומורפיזם חד־חד ערכי ועל.
	יהיו $\frac{r}{s} / \frac{t}{1} = \frac{r'}{s'} / \frac{t'}{1}$, כלומר קיים $a \in T$ כך ש־$\frac{r}{s} t' a = \frac{r'}{s'} t a$.
	לכן גם $\frac{r t' a}{s} = \frac{r' t a}{s'}$ וקיים $b \in S$ כך ש־$r t' a s' b = r' t a s b$, ולכן על־ידי בחירת $a b \in S T$ נקבל בדיוק ש־$\frac{r}{s t} = \frac{r'}{s' t'}$, כלומר ההעתקה מוגדרת היטב.
	נרצה להראות שזהו הומומורפיזם חוגים,
	\[
		\varphi(\frac{r}{s} / \frac{t}{1} + \frac{r'}{s'} / \frac{t'}{1})
		= \varphi(\frac{\frac{r}{s} t' + \frac{r'}{s'} t}{\frac{t t'}{1}})
		= \varphi(\frac{\frac{r t'}{s} + \frac{r' t}{s'}}{\frac{t t'}{1}})
		= \varphi(\frac{\frac{r t' s' + r' t s}{s s'}}{\frac{t t'}{1}})
		= \frac{r t' s' + r' t s}{s s' t t'}
		= \frac{r}{s t} + \frac{r'}{s' t'}
		= \varphi(\frac{r}{s} / \frac{t}{1}) + \varphi(\frac{r'}{s'} / \frac{t'}{1})
	\]
	ושימור כפל מתקבל באופן דומה (אך יותר פשוט), ונסיק שזהו אכן הומומורפיזם.
	חד־חד ערכיות נובעת באופן זהה מבדיקת הגדרה היטב, וגם על על־ידי בחירת נציגים, ולכן זהו איזומורפיזם.
\end{proof}

\question{}
יהי $I \triangleleft R$ ונגדיר את הקבוצה הכפלית $S = 1 + I$.
נראה כי $S^{-1} I \subset \Jj(S^{-1} R)$.
\begin{proof}
	נבחין כי כל איבר ב־$S^{-1} I$ הוא מהצורה $\frac{a}{1 + b}$ עבור $a, b \in I$, וכן שגם $\frac{a}{1 + b} \in S^{-1} R$.
	יהי $\frac{r}{1 + c} \in S^{-1} R$, ונבחן את $1 - \frac{a}{1 + b} \cdot \frac{r}{1 + c}$,
	\[
		1 - \frac{a}{1 + b} \cdot \frac{r}{1 + c}
		= \frac{1}{1} + \frac{- a r}{(1 + b)(1 + c)}
		= \frac{(1 + b)(1 + c) - a r}{(1 + b)(1 + c)}
	\]
	בהגדרה השקולה של רדיקל ג'ייקובסון מתקבל שאם ביטוי זה הפיך (עבור איברים כלליים) אז הוא שייך לרדיקל, ולכן נותר לבדוק אם לביטוי זה יש הופכי ב־$S^{-1} R$.
	ברור כי $(1 + b)(1 + c) \in R$, ולכן רק נרצה להראות ש־$(1 + b)(1 + c) - a r \in S$,
	\[
		(1 + b)(1 + c) - a r
		= 1 + b + c + b c - a r
		\in 1 + I
		\iff b + c + b c - a r
		\in I
	\]
	אבל הנחנו ש־$a, b, c \in I$, ולכן מסגירות לכפל גם $a r \in I$ ונסיק את הרצוי.
\end{proof}

\question{}
\subquestion{}
יהי $R$ תחום שלמות,
נראה שלכל $0 \notin S \subset R$ כפלית, גם $S^{-1} R$ תחום שלמות.
\begin{proof}
	נניח בשלילה ש־$S^{-1} R$ לא תחום שלמות, כלומר קיימים $\frac{r}{s}, \frac{r'}{s'} \in S^{-1} R$ כך ששניהם שונים מאפס אבל $\frac{r r'}{s s'} = 0$.
	יהי $a \in S$ כך שמתקיים $r r' s s' a = 0 a = 0$, אז נקבל ש־$r r' s s' \in R$ הוא מחלק אפס בסתירה.
\end{proof}

\subquestion{}
נראה שלכל $P \triangleleft \ZZ / 6 \ZZ$ ראשוני, ${(\ZZ / 6 \ZZ)}_P$ תחום שלמות.
\begin{proof}
	לחוג זה יש מספר סופי של קבוצות כפליות (של אידאלים ראשוניים), ואלו הם,
	\[
		\{ 1, 2, 4 \},
		\quad
		\{ 1, 3 \},
		\quad
		\{ 1, 5 \},
		\quad
		\{ 1, 3, 5 \},
		\quad
		\{ 1, 2, 4, 5 \}
	\]
	אלו בהכרח כל הקבוצות הכפליות כי $1 \in S$ תמיד מהגדרה וכן $0 \notin S$ כי $R \setminus S$ ראשוני ובפרט כולל איבר נייטרלי חיבורי.
	כדי לבדוק אם $S^{-1} R$ תחום שלמות, נוכל להראות ש־$r s \ne 0$ לכל $r \in R \setminus \{ 0 \}, s \in S$, ומבדיקה ישירה זה נכון.
\end{proof}

\question{}
נראה שקבוצה כפלית $0 \notin S \subset R$ היא רוויה אם ורק אם $R \setminus S$ איחוד של אידאלים ראשוניים.
\begin{proof}
	נאמר ש־$S$ רוויה אם $\forall x, y \in R, x y \in S \implies x, y \in S$.

	נניח כי $S$ רוויה, ויהי $a \in R \setminus S$.
	נבחן את $(a)$ האידאל הראשי הנוצר על־ידי $a$.
	אם $a^n \in S$ אז בהכרח גם $a \in S$ בסתירה, לכן $(a) \cap S = \varnothing$ בלבד.
	נוכל אם כך להרחיב את $(a)$ לאידאל מקסימלי זר ל־$S$ על־ידי משפט האידאל המקסימלי, נסמן $(a) \subset P_a \triangleleft R$ כך ש־$P_a \cap S = \varnothing$.
	בפרט $P_a$ הוא ראשוני, ולכן $R \setminus S = \bigcup_{a \in R \setminus S} P_a$ כפי שרצינו.

	בכיוון ההפוך נניח ש־$R \setminus S$ איחוד של אידאלים ראשוניים, אז ברור כי $S$ כפלית, ועלינו להראות כי היא רוויה.
	נניח בשלילה שהיא לא, ולכן קיימים $x, y \in R$ כך ש־$x y \in S$ אבל $x \notin S$, ובהתאם $x \in P$ עבור איזשהו אידאל ראשוני.
	לכן גם $x y \in P$ בסתירה.
\end{proof}

\question{}
תהי $S \subset R$ קבוצה כפלית, ונסמן $\bar{S}$ הריווי של $S$.
נראה שמתקיים,
\[
	S^{-1} R
	\cong \bar{S}^{-1} R
\]
\begin{proof}
	ראינו כי הקבוצה,
	\[
		\{ x \in R \mid \frac{x}{1} \in {(S^{-1} R)}^\times \}
		= \{ x \in R \mid \exists r \in R,\ r x \in S \}
	\]
	היא הריווי של $S$, כלומר הקבוצה המינימלית הכפלית הרוויה הכוללת את $S$.
	נגדיר את ההעתקה,
	\[
		\varphi : S^{-1} R \to \bar{S}^{-1} R,
		\qquad
		\varphi(\frac{r}{s}) = \frac{r}{s}
	\]
	ברור כי העתקה זו מוגדרת היטב, והומומורפיזם חוגים, וכן ברור שהיא חד־חד ערכית, נראה להראות על.
	יהי $s' \in \bar{S}$, כלומר $\frac{s'}{1} \in {(S^{-1} R)}^\times$, ולכן גם $\frac{1}{s'} \in S^{-1} R$.
	לכן אם $\frac{r}{s'} \in \bar{S}^{-1} R$ אז $\varphi(\frac{r}{s'})$ מוגדר ושולח בדיוק לאיבר המבוקש.
\end{proof}

\question{}
יהי $A$ תחום ראשי ויהי $P \triangleleft A$ ראשוני.
נראה ש־$A_P$ תחום ראשי.
\begin{proof}
	יהי $I \triangleleft A_P$ אידאל, ונראה שהוא ראשי.
	ראינו כי כל אידאל במיקום הוא מהצורה $S^{-1} J$ כאשר $S = A \setminus P$ ו־$J \triangleleft A$.
	אבל $A$ ראשי ולכן $J = (a)$ לאיזשהו $a \in A$, ובהתאם גם $I = S^{-1} (a) = \{ \frac{x}{s} \mid x \in (a), s \in S \}$.
	אבל לכל $\frac{x}{s}$ כזה נבחין כי $\frac{x}{s} = \frac{x}{1} \cdot \frac{1}{s}$ כאשר $\frac{1}{s} \in A_P$ וכן $\frac{x}{1} = \loc(x)$ עבור $x \in (a)$, ולכן נסיק שב־$A_P$ מתקיים $(a) = I$ בדיוק.
\end{proof}

\question{}
יהי $A$ חוג ויהי $f \in A$, עליו נגדיר את הקבוצה הכפלית $S = \{ f^n \mid n \in \NN \}$, ונסמן את $A_f = S^{-1} A$.
נוכיח כי,
\[
	A_f
	\cong A[x] / (1 - x f)
\]
\begin{proof}
	נגדיר את,
	\[
		\varphi(\frac{a}{f^n})
		= a x^n + (1 - x f)
	\]
	ונבדוק שהעתקה זו מוגדרת היטב, נניח ש־$\frac{a}{f^n} = \frac{b}{f^m}$ אז קיים $k \in \NN$ כך ש־$a f^m f^k = b f^n f^k$ וכן,
	\[
		\varphi(\frac{a}{f^n})
		= a x^n + (1 - x f)
		= a x^n - a x^n (1 - x f) + (1 - x f)
		= a x^n \cdot x f + (1 - x f)
	\]
	כלומר $t \equiv t \cdot x f \mod (1 - x f)$ ולכן,
	\begin{multline*}
		\varphi(\frac{a}{f^n})
		= a x^n + (1 - x f)
		= a x^n \cdot x^{m + k} f^{m + k} + (1 - x f)
		= b x^n \cdot x^{m + k} f^{n + k} + (1 - x f) \\
		= b x^m \cdot x^{n + k} f^{n + k} + (1 - x f)
		= b x^m + (1 - x f)
		= \varphi(\frac{b}{f^m})
	\end{multline*}
	ומצאנו שההעתקה אכן מוגדרת היטב.
	ברור כי היא חד־חד ערכית ועל, וכן כי זהו הומומורפיזם מבדיקה ישירה.
\end{proof}

\question{}
יהי $f : A \to B$ הומומורפיזם חוגים, ותהי $S \subset A$ כפלית.
נסמן את $T = f(S)$, זוהי קבוצה כפלית, ובהתאם $T^{-1} B$ הוא מודול מעל $S^{-1} A$.
$B$ הוא גם מודול מעל $A$ ולכן נוכל להגדיר גם את המיקום של $B$ ב־$S$, כלומר $S^{-1} B$ הוא מודול מעל $S^{-1} A$.
נראה כי,
\[
	S^{-1} B
	\cong T^{-1} B
\]
כמודולים מעל $S^{-1} A$.
\begin{proof}
	נרצה לבנות את המיפוי $\frac{b}{s} \mapsto \frac{b}{f(s)}$ ולהראות שהוא איזומורפיזם $S^{-1} B \to T^{-1} B$, אך זה ברור.
\end{proof}

\question{}
יהי $A$ תחום שלמות ויהי $F = \Frac(A)$ שדה השברים ל $A$, כלומר $F = {(A \setminus \{ 0 \})}^{-1} A$, ונזהה את $a \in A$ עם $\frac{a}{1} \in \Frac(A)$ לכל $a$. \\
לכל אידאל מקסימלי $M \triangleleft A$, $A_M$ הוא תת־חוג של $F$.
נראה כי,
\[
	A
	= \bigcap_{\substack{M \triangleleft A \\ M \text{ maximal}}} A_M
\]
\begin{proof}
	לכל $M$ נבחין כי $\frac{a}{1} \in A_M$ לכל $a \in A$, ולכן $A \subset A_M$, בפרט,
	\[
		A
		\subset \bigcap_M A_M
	\]
	ונותר להראות את הכיוון ההפוך.

	יהי $x \in \bigcap_M A_M$, אז נגדיר את $I = \{ a \in A \mid a x \in A \}$, אז $I \triangleleft A$.
	אם $x \in A$ אז סיימנו, ולכן נניח ש־$x \notin A$ ולכן גם $1 \notin I$, כלומר $I$ הוא אידאל ממש, ובהתאם קיים אידאל מקסימלי $I \subset M \triangleleft A$.
	מתקיים $x \in A_M$ ולכן $x = \frac{a}{s}$ עבור $a \in A$ ו־$s \in A \setminus M$, וכן $s \cdot x = s \cdot \frac{a}{s} = \frac{a}{1} = a$ ולכן $s \in I$, ובפרט $s \in M$ בסתירה ל־$s \notin M$, לכן $x \in A$ בלבד.
\end{proof}

\question{}
\subquestion{}
יהי $R$ חוג, $M$ מודול נוצר סופית מעל $R$, ותהי $S \subset R$ כפלית.
נראה ש־$S^{-1} M = 0$ אם ורק אם קיים $s \in S$ כך ש־$s M = 0$.
\begin{proof}
	נניח ש־$S^{-1} M = 0$, כלומר לכל $\frac{m}{s} \in S^{-1} M$ קיים $s_m \in S$ כך ש־$m s_m = 0$.
	נסמן $e_1, \ldots, e_n \in M$ קבוצה פורשת של $M$, ונסמן גם $s_1, \ldots, s_n$ כך ש־$s_i = s_{m_i}$.
	לכל $m \in M$ מתקיים $m = \sum_i r_i e_i$ ולכן,
	\[
		(e_1 \cdots e_n) \cdot m
		= \sum_{i = 1}^n (s_1 \cdots s_n) \alpha_i e_i
		= \sum_{i = 1}^n (s_1 \cdots s_{i - 1} \cdot s_{i + 1} \cdots s_n) \alpha_i s_i e_i
		= \sum_{i = 1}^n (s_1 \cdots s_{i - 1} \cdot s_{i + 1} \cdots s_n) \alpha_i \cdot 0
		= 0
	\]
	ולכן $s = e_1 \cdots e_n$ מקיים $s M = 0$.

	בכיוון ההפוך נניח ש־$s M = 0$ עבור $s \in S$.
	לכן לכל $\frac{m}{t} \in S^{-1} M$ מתקיים $m \cdot 1 \cdot s = 0 \cdot t \cdot s$, כלומר $\frac{m}{t} = \frac{0}{1}$, ונסיק ש־$S^{-1} M = 0$.
\end{proof}

\subquestion{}
נמצא חוג $R$ ומודול $M$ מעליו שאינו נוצר סופית, וקבוצה כפלית $S \subset R$ כך ש־$S^{-1} M = 0$ אבל לא קיים $s \in S$ כך ש־$s M = 0$.
\begin{solution}
	נגדיר את $R = \ZZ$ ואת $M = \bigoplus_{n \in \NN} \ZZ / n \ZZ$, וכן $S = \ZZ \setminus \{ 0 \}$, זוהי קבוצה כפלית.
	יהי $\frac{(a_1, a_2, \ldots)}{b} \in S^{-1} M$, אז $a_i = 0$ למעט מספר סופי של אינדקסים, לכן,
	\[
		s
		= \prod_{\substack{i \in \NN \\ a_i \ne 0}} i
	\]
	ולכן $s \cdot (a_1, \ldots) = (s a_1, s a_2, \ldots) = (0, 0, \ldots)$ בלבד, כלומר $\frac{(a_i)}{b} = 0$ ולכן $S^{-1} M = 0$.

	נניח בשלילה שקיים $s \in S$ כך שגם $s M = 0$, ונבחר $e_i \in M$ מהבסיס הסטנדרטי, אז $s e_i = 0 \implies s \equiv 0 \mod i$, כלומר $i \mid m$, אבל זה נכון לכל $i \in \NN$ בסתירה ל־$s \in \ZZ \setminus \{ 0 \}$.
\end{solution}

\question{}
יהי $A$ חוג, $M$ מודול מעליו, ו־$I \triangleleft A$.
נניח שלכל $I \subset N \triangleleft A$ מקסימלי מתקיים $M_N = 0$.
נראה ש־$M = I M$.
\begin{proof}
	הכיוון של $I M \subset M$ ברור. \\
	נגדיר את $C = M / I M$, ואנו רוצים להראות ש־$C \cong \{ I M \}$, ולכן נניח בשלילה ש־$c = m + I M \in C$ איבר לא טריוויאלי.
	מהנתון ומשימור מיקום מתקיים $C_N = M_N / {(I M)}^N = M_N / I_N M_N = 0$ לכל $I \subset N$ אידאל מקסימלי.

	יהי $N$ אידאל מקסימלי כלשהו. אם $I \subset N$ אז $C_N = 0$, נרצה להראות שגם כאשר $I \not\subset N$ זהו המצב.
	במצב זה $N + I = A$ בלבד, וקיים $s \in I \setminus N$, לכן $s \in R \setminus N$ ונובע ש־$\frac{s}{1}$ הפיך ב־$C_N$.
	נסיק שגם $\frac{s}{1} \in I_N$, אבל גם $\frac{1}{s} \in I_N$ ולכן $I_N = A_N$ ונסיק שגם $I_N M_N = A_N M_N = M_N$ ולכן $C = M_N / M_N = \{ I M \}$, וקיבלנו שתמיד $C$ טריוויאלי.
\end{proof}

\question{}
יהי $A$ תחום שלמות.
$A$־מודול מעל $M$ נקרא חסר פיתול אם $T(M) = \{ 0 \}$.
נראה שלכל $S \subset A$ כפלית ו־$M$ מודול מתקיים,
\[
	T(S^{-1} M)
	= S^{-1} T(M)
\]
ונסיק ש־$M$ חסר פיתול אם ורק אם לכל $P \triangleleft A$ ראשוני, $M_P$ חסר פיתול.
\begin{proof}
	ניזכר ש־$T(M) = \{ m \in M \mid \Ann(m) \ne \{ 0 \} \}$.
	יהי $\frac{m}{s} \in T(S^{-1} M)$, ונרצה להראות ש־$m \in T(M)$.
	מהגדרת המאפס נסיק שקיים איבר לא טריוויאלי $a \in A$ כך ש־$a \frac{m}{s} = 0$, כלומר $a t \cdot m = 0$ לאיזשהו $t \in S$, אבל נובע אם כך ש־$a t \in \Ann(m)$, כלומר $m \in T(M)$ כפי שרצינו.

	מהצד השני יהי $\frac{m}{s} \in S^{-1} T(M)$, ונרצה להראות שגם $\frac{m}{s} \in T(S^{-1} M)$.
	מהנתון קיים $a \in A$ כך ש־$a m = 0$, אז בפרט גם $\frac{a m}{s} = \frac{0}{1}$ וקיבלנו בדיוק את המבוקש.

	ניזכר כי $M = 0$ היא תכונה מקומית, לכן הטענה נכונה גם ל־$T(M)$, כלומר $T(M) = 0$ (חסרת פיתול) אם ורק אם לכל $P \triangleleft A$ ראשוני, ${(T(M))}_P = T(M_P) = 0$.
\end{proof}

\question{}
יהי חוג $R$, ונגדיר,
\[
	\dim R
	= \sup\{ n \in \NN \mid \exists P_0 \subsetneq P_1 \subsetneq \cdots \subsetneq P_n \text{ prime ideals of } R \}
\]
ונגדיר את מימד Krull של $R$ להיות $\dim R$.

\subquestion{}
נניח ש־$R$ תחום שלמות, ונראה ש־$R$ שדה אם ורק אם $\dim R = 0$.
\begin{proof}
	אם $R$ שדה אז כל אידאל של $R$ הוא טריוויאלי, ולכן השרשרת הארוכה ביותר של אידאלים ראשוניים היא $\{ 0 \} \triangleleft R$, כלומר $\dim R = 0$.

	מהצד השני נניח ש־$\dim R = 0$, ונסמן את השרשרת המעידה על־כך $P_0 \triangleleft R$.
	אם $P_0$ לא טריוויאלי אז נוכל לבחור $\{ 0 \} \subsetneq P_0 \triangleleft R$, לכן $\dim R \ge 1$ בסתירה, ונסיק ש־$P_0 = \{ 0 \}$ בלבד.
	נובע של־$R$ אין אידאלים ראשוניים לא טריוויאליים, ולכן בפרט הטענה נכונה גם עבור אידאלים מקסימליים, כלומר $\{ 0 \}$ אידאל מקסימלי יחיד, ונובע ש־$R = R / \{ 0 \}$ שדה.
\end{proof}

\subquestion{}
יהי $P \triangleleft R$ ראשוני, נראה כי,
\[
	\dim R_P
	= \sup\{ n \in \NN \mid \exists P_0 \subsetneq P_1 \subsetneq \cdots \subsetneq P_n = P \text{ prime ideals of } R \}
\]
\begin{proof}
	נניח כי $\dim R_P = n$ ותהי $Q_0 \subsetneq \cdots Q_n$ שרשרת המעידה על כך, כלומר $Q_i \triangleleft R_P$ ראשוני לכל $i \le n$.
	נסמן $S = R \setminus P$ ויהיו $Q_i = S^{-1} P_i$ לכל $i \le n$ כאשר $P_i \triangleleft R$ אידאל ראשוני. \\
	נבחין כי לכל $0 < i \le n$,
	\[
		x \in P_i \setminus P
		\implies x \in S
		\implies \frac{1}{x} \in R_P
		\implies \frac{x}{x} \in P_i
	\]
	וזו כמובן סתירה ולכן $P_i \subset P$ לכל $i$ (גם ל־$0$), וממקסימליות השרשרת נסיק ש־$P_n = P$ בלבד, אחרת נוכל לסמן $P_{n + 1} = P$ ולקבל $\dim R_P = n + 1$.
\end{proof}

\subquestion{}
נחשב את $\dim \ZZ$, וכן לכל $p \in \ZZ$ ראשוני, נחשב את $\dim \ZZ_{(p)}$.
\begin{solution}
	ניזכר ש־$P \triangleleft \ZZ$ הוא ראשוני אם ורק אם $P = (p)$ לאיזשהו ראשוני $p$ או $P = (0)$, ונסמן $P_p = (p)$ לכלל האידאלים של $\ZZ$.
	בהתאם נובע ש־$P_p \not\subset P_q$ לכל $p, q$ ראשוניים, ובפרט $(0) \subsetneq P_2 \triangleleft \ZZ$ היא שרשרת מקסימלית, לכן $\dim \ZZ = 1$.

	מהסעיף הקודם נובע ש־$(0) \subsetneq P_p$ שרשרת המעידה על $\dim \ZZ_{(p)} = 1$.
\end{solution}

\end{document}
